\documentclass[11pt,a4paper]{article}
\usepackage[utf8]{inputenc}
\usepackage[T1]{fontenc}
\usepackage[spanish,es-nodecimaldot,es-noshorthands]{babel}
\usepackage{lmodern}
\usepackage{amsmath,amssymb}
\usepackage[margin=2.4cm,headheight=14pt]{geometry}
\usepackage{booktabs,tabularx,array}
\usepackage{enumitem}
\usepackage{xcolor}
\usepackage{graphicx}
\usepackage{listings}
\usepackage{fancyhdr}
\usepackage[strings]{underscore}
\usepackage{hyperref}

\definecolor{unl}{RGB}{0,86,145}
\definecolor{gris}{HTML}{475569}
\hypersetup{colorlinks=true,linkcolor=unl,urlcolor=unl}
\setlength{\parskip}{5pt}
\setlength{\parindent}{0pt}
\setlist{leftmargin=*,itemsep=2pt,topsep=3pt}
\newcolumntype{Y}{>{\raggedright\arraybackslash}X}
\newcolumntype{L}[1]{>{\raggedright\arraybackslash}p{#1}}
\lstset{basicstyle=\ttfamily\footnotesize,breaklines=true,frame=single,rulecolor=\color{gris},
  columns=fullflexible,keepspaces=true,showstringspaces=false,
  literate={á}{{\'a}}1 {é}{{\'e}}1 {í}{{\'i}}1 {ó}{{\'o}}1 {ú}{{\'u}}1 {ñ}{{\~n}}1
    {Á}{{\'A}}1 {É}{{\'E}}1 {Í}{{\'I}}1 {Ó}{{\'O}}1 {Ú}{{\'U}}1 {Ñ}{{\~N}}1 {ü}{{\"u}}1}
\pagestyle{fancy}
\fancyhf{}
\lhead{Matemática II: Cálculo Integral}
\rhead{Avance 2 -- Proyecto integrador}
\cfoot{\thepage}
\renewcommand{\headrulewidth}{0.4pt}
% Texto gris = instrucción. Reemplácelo o elimínelo antes de entregar.
\newcommand{\completar}[1]{\textit{\textcolor{gris}{#1}}}
\newcommand{\completarm}[1]{{\color{gris}#1}}

\begin{document}
\hypersetup{pageanchor=false}

\begin{titlepage}
\centering
\vspace*{1.6cm}
{\Large\bfseries UNIVERSIDAD NACIONAL DE LOJA\par}
\vspace{0.3cm}
{\large Facultad Agropecuaria y de Recursos Naturales Renovables\par}
{\large Carrera de Ingeniería Ambiental\par}
\vspace{1.2cm}
{\color{unl}\rule{0.82\linewidth}{1.3pt}\par}
\vspace{0.5cm}
{\huge\bfseries \completar{Título: magnitud integral, tema y zona}\par}
\vspace{0.4cm}
{\Large Proyecto integrador -- Avance 2\par}
{\large Unidad 1, Fase 2: Avance técnico y datos\par}
\vspace{0.5cm}
{\color{unl}\rule{0.82\linewidth}{1.3pt}\par}
\vfill
\begin{tabular}{rl}
\textbf{Asignatura:} & Matemática II: Cálculo Integral\\
\textbf{Grupo:} & \completar{Número de grupo según la nómina oficial}\\
\textbf{Tema y zona:} & \completar{Bloque temático -- Norte, Sur o cantón completo}\\
\textbf{Integrantes:} & \completar{Apellidos y nombres completos}\\
 & \completar{Añada una fila por integrante}\\
\textbf{Docente:} & Ing. Guillermo Chuncho\\
\textbf{Período académico:} & septiembre 2026 -- febrero 2027\\
\textbf{Fecha:} & \completar{dd de mes de 2026}\\
\end{tabular}
\vspace{1cm}
\end{titlepage}
\hypersetup{pageanchor=true}

% =====================================================================
\section{Respuesta a la revisión del avance 1}
Respondan a cada observación de su hoja en \texttt{Matriz\_revision\_avance1\_matematica.xlsx}. Incluyan obligatoriamente todas las observaciones críticas y mayores. Si no están de acuerdo con alguna, argumenten con evidencia. Los grupos sin entrega previa escriben «Primera presentación».

\small
\begin{tabularx}{\linewidth}{L{1.0cm} Y Y L{2.0cm}}
\toprule
\textbf{N.º} & \textbf{Observación (resumen)} & \textbf{Cambio realizado o justificación} & \textbf{Sección / página}\\
\midrule
\completar{1} & \completar{Resumen de la observación} & \completar{Qué se cambió y por qué} & \completar{Sec. x}\\
\completar{2} & \completar{…} & \completar{…} & \completar{…}\\
\bottomrule
\end{tabularx}
\normalsize

% =====================================================================
\section{Problema de investigación}
\subsection{Planteamiento del problema}
Organicen el texto en este orden: (1) situación territorial o ambiental de la ventana de estudio; (2) vacío de conocimiento: qué magnitud acumulada no se ha cuantificado; (3) relevancia, sin atribuir a la base efectos que no mide; (4) cómo una integral definida y los datos asignados permiten resolverlo.

\completar{Redacte aquí el problema con datos de contexto citados (APA 7). Elimine estas instrucciones antes de entregar.}

\subsection{Pregunta de investigación}
Una sola pregunta, en clave de cálculo integral, que nombre la magnitud y sus unidades, la zona (o el transecto) y el periodo, y que pueda responderse con el CSV.

\textbf{Pregunta:} \completar{Escriba aquí la pregunta.}

\subsection{Hipótesis}
Una afirmación contrastable, con indicador y umbral cuantitativo justificados. Escriban por separado el criterio de validación numérica del modelo.

\textbf{Hipótesis:} \completar{Afirmación medible sobre el fenómeno.}

\textbf{Justificación del umbral:} \completar{Fuente bibliográfica o evidencia de la exploración preliminar.}

\textbf{Criterio de validación del modelo:} \completar{Ej.: error relativo entre la integral del modelo y la suma discreta menor que X \%.}

% =====================================================================
\section{Objetivos}
\subsection{Objetivo general}
\textbf{Objetivo general:} \completar{Verbo en infinitivo + magnitud integral + zona + datos + periodo. Debe responder la pregunta.}

\subsection{Objetivos específicos}
\begin{enumerate}
  \item \completar{Objetivo específico 1.} \textbf{Salida:} \completar{tabla, valor, figura o error \%.}
  \item \completar{Objetivo específico 2.} \textbf{Salida:} \completar{…}
  \item \completar{Objetivo específico 3.} \textbf{Salida:} \completar{…}
\end{enumerate}

% =====================================================================
\section{Metodología}
\subsection{Zona de estudio y fuente de datos}
\textbf{Grupo y tema:} \completar{Número de grupo y bloque temático.}\\
\textbf{Zona y regla territorial:} \completar{Norte o Sur según centro\_y\_m respecto de 9\,555\,250 m, o cantón completo; indique cómo lo verificaron.}\\
\textbf{Ventana de estudio:} \completar{Filas y columnas de la malla; número de celdas; área en km$^2$.}\\
\textbf{Archivo asignado:} \completar{Nombre exacto y ruta en Unidad\_1/base\_de\_datos.}\\
\textbf{Periodo y unidad de observación:} \completar{2019--2024; celda de 500 m (relieve) o celda--mes (cobertura, clima), o mes cantonal (incendios).}

\subsection{Marco de variables}
Usen los nombres exactos de las columnas y consulten el significado y la unidad en \texttt{DICCIONARIO\_VARIABLES.csv}. Incluyan solo las variables que usarán y distingan las analíticas de las de control.

\small
\begin{tabularx}{\linewidth}{L{3.4cm} Y L{1.0cm} L{0.9cm} L{1.5cm} L{1.3cm} L{1.6cm}}
\toprule
\textbf{Columna} & \textbf{Magnitud} & \textbf{Símb.} & \textbf{Dim.} & \textbf{Unidad dato} & \textbf{Unidad SI} & \textbf{Uso}\\
\midrule
\completar{nombre\_exacto} & \completar{Definición} & \completar{$x$} & \completar{L} & \completar{m} & \completar{m} & \completar{analítica / control}\\
\completar{nombre\_exacto} & \completar{Definición} & \completar{} & \completar{} & \completar{} & \completar{} & \completar{}\\
\completar{nombre\_exacto} & \completar{Definición} & \completar{} & \completar{} & \completar{} & \completar{} & \completar{}\\
\bottomrule
\end{tabularx}
\normalsize

\subsection{Auditoría de calidad}
Reporten el resultado de cada control (realizado con el código del anexo).

\small
\begin{tabularx}{\linewidth}{L{4.6cm} L{3.0cm} Y}
\toprule
\textbf{Control} & \textbf{Esperado} & \textbf{Observado / decisión}\\
\midrule
Filas $\times$ columnas & \completar{p. ej. 7\,200 $\times$ 22} & \completar{}\\
Celdas únicas & \completar{} & \completar{}\\
Periodo (anio, mes, t\_mes); sin 2025 & 2019--2024; 1--12; 1--72 & \completar{}\\
Regla de zona (centro\_y\_m) & \completar{} & \completar{}\\
Duplicados & 0 & \completar{}\\
Faltantes por variable analítica & \completar{} & \completar{n y \% por columna; por mes o por celda si aplica}\\
Campos de calidad del tema & \completar{p. ej. sentinel\_t1\_missing} & \completar{}\\
Rangos físicamente plausibles & \completar{} & \completar{}\\
\bottomrule
\end{tabularx}
\normalsize

\textbf{Tratamiento de datos faltantes:} \completar{Excluir, promediar las celdas válidas o interpolar; justifique la decisión.}

\subsection{Exploración preliminar}
\begin{itemize}
  \item \textbf{Cinco filas de ejemplo:} \completar{tabla breve con las columnas seleccionadas; indique cómo se eligieron.}
  \item \textbf{Estadísticos descriptivos:} \completar{tabla con $n$ válido, mínimo, máximo, media y desviación estándar, con unidades.}
  \item \textbf{Primera visualización exploratoria:} \completar{inserte la figura con título, ejes, unidades, zona y periodo.}
  \item \textbf{Hechos observados:}
  \begin{enumerate}
    \item \completar{Hecho 1 y la variable o los registros que lo sustentan.}
    \item \completar{Hecho 2 …}
    \item \completar{Hecho 3 …}
  \end{enumerate}
\end{itemize}

% Ejemplo para insertar la figura generada por el código:
% \begin{figure}[h]\centering
%   \includegraphics[width=0.85\linewidth]{figuras/exploracion_1.pdf}
%   \caption{Título, zona y periodo. Fuente: elaboración propia con base en ...}
% \end{figure}

% =====================================================================
\section{Planteamiento matemático inicial}
\subsection{Modelo funcional}
\completar{Variable independiente, dominio y unidades; forma propuesta de la función (polinómica, trigonométrica, por tramos…) y por qué, según la exploración.}
\[
  \completarm{f(x) =\ldots \qquad x\in[a,b]\ \text{[unidades]}}
\]

\subsection{Integral definida}
\completar{Integral que se calculará, límites de integración justificados, unidades del resultado e interpretación física.}
\[
  \completarm{I =\int_{a}^{b} f(x)\,dx \quad [\text{unidades}]}
\]
\textbf{Conexión con el TFC:} \completar{Antiderivada $F$, qué representa y cómo se verificará que $\int_a^b f = F(b)-F(a)$.}

\subsection{Discretización y aproximación inicial}
\completar{$\Delta x$ o $\Delta t$, número de subintervalos $n$ y fórmula de la suma de Riemann o del trapecio con los datos observados.}
\[
  \completarm{T_n =\frac{\Delta x}{2}\left[f(x_0) + 2\sum_{i=1}^{n-1} f(x_i) + f(x_n)\right]}
\]
\textbf{Resultado preliminar:} \completar{valor obtenido con el código, con unidades. Indique que es preliminar.}

\textbf{Control dimensional:} \completar{p. ej. [m$^2$] = [m] $\times$ [m].}

\subsection{Actividades realizadas y pendientes}
\completar{Separe lo que ya se comprobó de lo que queda planificado para el producto técnico (Fase 3).}

% =====================================================================
\section*{Referencias}
\completar{Lista en APA 7, solo con fuentes citadas en el texto y localizables (DOI o URL). Incluya la base y el diccionario del curso.}

\begin{itemize}
  \item \completar{Autor, A. (Año). Título. \textit{Revista, vol}(n), pp. https://doi.org/...}
\end{itemize}

\appendix
\section{Código reproducible}
\textbf{Ruta en el repositorio:} \completar{p. ej. Unidad\_1/proyecto/grupo\_XX/avance2.ipynb}

\begin{lstlisting}[language=Python]
# Fragmento representativo (reemplazar por el código del equipo)
import pandas as pd
df = pd.read_csv("ruta/relativa/al/archivo.csv")   # solo lectura
print(df.shape)
\end{lstlisting}

\section{Declaración de uso de inteligencia artificial}
\completar{Herramienta utilizada, para qué se usó, qué verificó el equipo y qué partes fueron reescritas por el equipo.}

\end{document}
